% ---------------------------------------------------------------------
% 2.6 Procesamiento digital de señales biomédicas
% ---------------------------------------------------------------------
\section[PROCESAMIENTO DIGITAL DE SEÑALES BIOMÉDICAS]{Procesamiento Digital de Señales
Biomédicas}
\label{sec:pds}

El preprocesamiento de las señales EEG constituye una etapa determinante para el desempeño de
los modelos de clasificación, dado que estas señales presentan una relación señal-ruido baja y
son susceptibles a artefactos de origen fisiológico y ambiental. Para atender esta problemática
se emplean técnicas de filtrado pasa banda, orientadas a delimitar el rango de frecuencias
asociado a la componente P300, así como la segmentación de la señal continua en ventanas
temporales alrededor de cada estímulo y la normalización de las mismas, cuyos fundamentos se
describen a lo largo de esta sección.

\subsection{Muestreo y teorema de Nyquist}

Las señales fisiológicas son analógicas, es decir, están definidas en todo instante y su amplitud
puede tomar cualquier valor. Para procesarlas en una computadora se deben muestrear a intervalos
regulares $T_s = 1/f_s$, lo que produce la secuencia discreta $x(n) = x_a(nT_s)$
\cite{proakis2007}. El teorema
de muestreo establece que una señal cuyo contenido en frecuencia está limitado a $f_{\max}$ puede
reconstruirse sin pérdida a partir de sus muestras siempre que
\begin{equation}
f_s \geq 2 f_{\max},
\label{eq:nyquist}
\end{equation}
y si esta condición no se cumple, las componentes de frecuencia superiores a $f_s/2$, conocida
como frecuencia de Nyquist, se reflejan sobre las frecuencias bajas en un fenómeno llamado
\textit{aliasing} \cite{oppenheim2010}. La base de datos utilizada en este trabajo fue muestreada a
240~Hz, por lo que su frecuencia de Nyquist de 120~Hz está muy por encima de la banda de interés
de la P300, que se concentra por debajo de 20~Hz.

\subsection{Transformada de Fourier y densidad espectral de potencia}

La transformada discreta de Fourier (DFT) permite representar una señal como una suma de
componentes sinusoidales de distintas frecuencias, y en la práctica se calcula con el algoritmo de
la transformada rápida de Fourier (FFT) \cite{proakis2007}. A partir de ella se estima la densidad
espectral de potencia (PSD), que describe cómo se distribuye la potencia de la señal en función de
la frecuencia. En el caso del EEG suele emplearse el método de Welch, que divide la señal en
segmentos traslapados, calcula el espectro de cada uno y los promedia para obtener una estimación
más estable \cite{oppenheim2010}. Contreras utilizó este análisis para identificar que la
información útil para detectar la P300 en sus registros se encontraba entre 1 y 25~Hz
\cite{contreras2021}. La Figura~\ref{fig:psd} muestra una PSD típica, en la que se observan el pico
del ritmo alfa y la interferencia de la red eléctrica.

\begin{figure}[htbp]
\centering
\includegraphics[width=0.8\textwidth]{02-marco-referencial/psd_welch.pdf}
\caption{Densidad espectral de potencia estimada con el método de Welch para una señal de EEG
sintética. Elaboración propia.}
\label{fig:psd}
\end{figure}

\subsection{Convolución y sistemas lineales}

Un sistema lineal e invariante en el tiempo queda completamente descrito por su respuesta al
impulso $h(n)$, y su salida ante cualquier entrada $x(n)$ se obtiene mediante la operación de
convolución \cite{proakis2007}:
\begin{equation}
y(n) = (x * h)(n) = \sum_{k=-\infty}^{\infty} x(k)\, h(n-k).
\label{eq:convolucion}
\end{equation}
Esta operación es la base de los filtros digitales que se describen a continuación y, como se
verá en la Sección~\ref{sec:cnn}, también de las capas de las redes neuronales convolucionales,
en las que los coeficientes de $h(n)$ ya no se diseñan manualmente sino que se aprenden a partir
de los datos \cite{goodfellow2016}.

\subsection{Filtros digitales}

Los filtros digitales se clasifican en dos grandes familias. Los filtros de respuesta finita al
impulso (FIR) calculan la salida solo a partir de las muestras de entrada y siempre son estables.
Los filtros de respuesta infinita al impulso (IIR) también usan salidas anteriores, lo que les
permite lograr una buena selectividad con un orden mucho menor \cite{proakis2007,oppenheim2010}.
Entre los filtros IIR, el de Butterworth es uno de los más utilizados en EEG porque su respuesta
en magnitud es máximamente plana en la banda de paso \cite{oppenheim2010}. De acuerdo con la banda
que dejan pasar, los filtros pueden ser pasa bajas, pasa altas, pasa banda o rechaza banda. El
filtro \textit{notch} es un tipo de filtro rechaza banda que elimina una sola frecuencia, como la
de la red eléctrica. La Figura~\ref{fig:filtros} muestra la respuesta en frecuencia de un filtro
pasa banda entre 0.1 y 20~Hz, como el aplicado en el trabajo preliminar, y de un filtro
\textit{notch} en 60~Hz.

\begin{figure}[htbp]
\centering
\includegraphics[width=\textwidth]{02-marco-referencial/respuesta_filtros.pdf}
\caption{Respuesta en magnitud de (a) un filtro pasa banda Butterworth de cuarto orden entre 0.1 y
20~Hz y (b) un filtro \textit{notch} en 60~Hz, para $f_s = 240$~Hz. Elaboración propia.}
\label{fig:filtros}
\end{figure}

\subsection{Remuestreo}

El remuestreo permite reducir la frecuencia de muestreo después de filtrar la señal. En la
decimación por un factor $M$ se conserva una de cada $M$ muestras. Antes de hacerlo, se debe
aplicar un filtro antialiasing que elimine el contenido por encima de la nueva frecuencia de
Nyquist \cite{proakis2007}. En el trabajo preliminar se redujo la frecuencia de 240~Hz a 120~Hz,
lo que disminuye a la mitad el tamaño de cada época sin perder información en la banda de 0.1 a
20~Hz.

\subsection{Referencia y corrección de línea base}

Además del filtrado en frecuencia, es común modificar la referencia de los canales después del
registro. La referencia promedio común resta a cada canal $x_{ch}(n)$ el promedio de los $C$
canales en cada instante,
\begin{equation}
x^{\text{CAR}}_{ch}(n) = x_{ch}(n) - \frac{1}{C}\sum_{c=1}^{C} x_{c}(n),
\label{eq:car}
\end{equation}
con lo que se atenúan las componentes comunes a todos los electrodos, como las derivadas de la
propia referencia \cite{luck2014}. De forma complementaria, la corrección de línea base resta a
cada época el valor promedio de un intervalo previo al estímulo, típicamente de 100 a 200~ms, para
que las amplitudes de los componentes del ERP se midan respecto a un nivel común \cite{luck2014}.

\subsection{Segmentación en épocas}

La segmentación consiste en dividir la señal continua en ventanas sincronizadas con cada estímulo,
que constituyen las muestras de entrada al clasificador. La longitud de la ventana debe ser
suficiente para contener la P300 completa. Contreras empleó ventanas de 1005~ms a partir del
estímulo \cite{contreras2021}, mientras que en el trabajo preliminar se utilizaron ventanas de 78
muestras, equivalentes a unos 650~ms a 120~Hz. Dado que el intervalo entre inicios de estímulos es
de solo 175~ms, las épocas de intensificaciones consecutivas se traslapan, como se observa en la
Figura~\ref{fig:epocas}, de modo que una época no objetivo puede contener parte de la P300
provocada por una intensificación objetivo cercana.

\begin{figure}[htbp]
\centering
\includegraphics[width=\textwidth]{02-marco-referencial/segmentacion_epocas.pdf}
\caption{Segmentación de la señal continua en épocas de 0 a 650~ms a partir de cada
intensificación. Señal sintética, elaboración propia.}
\label{fig:epocas}
\end{figure}

\subsection{Normalización}

Las amplitudes del EEG varían entre canales, sesiones y sujetos, por lo que antes de entrenar un
modelo se acostumbra normalizar los datos. La normalización \textit{z-score} transforma cada canal
para que tenga media cero y desviación estándar unitaria, mientras que la normalización
\textit{min-max}, utilizada por Contreras \cite{contreras2021}, escala los valores al intervalo
$[0,1]$:
\begin{equation}
z(n) = \frac{x(n)-\mu}{\sigma}, \qquad
x'(n) = \frac{x(n)-\min(x)}{\max(x)-\min(x)}.
\label{eq:normalizacion}
\end{equation}
La primera es menos sensible a valores atípicos, lo cual resulta conveniente en presencia de
artefactos de gran amplitud, y fue la empleada de manera independiente por canal en el trabajo
preliminar para evitar que los electrodos con mayor amplitud dominen el aprendizaje de la red.

\subsection{Remoción de artefactos}

Cuando los artefactos se traslapan en frecuencia con la señal de interés, como ocurre con los
parpadeos y la banda de la P300, el filtrado no basta para eliminarlos. La estrategia más simple
es descartar las épocas cuya amplitud supere un umbral, aunque así se pierden datos. Otra opción es
el análisis de componentes independientes (ICA), que separa la señal en fuentes independientes
para identificar y eliminar las asociadas a movimientos oculares o musculares \cite{jung2000}.
